\section{Finite BCH-256 certificate}
\label{app:finite-bch-proof}
\label{app:structured-finite-member}

This appendix connects the fixed encoder to the numerical inequalities in
Theorem~\ref{thm:finite-bch-spin}. The argument has three interfaces: a
deterministic outer spectrum bound, a positive inner transfer, and a complete
sum over outer occupancies. Computational witnesses supply finite upper
bounds at these interfaces; they do not replace the probability comparisons.

\subsection{The fixed outer and its spectrum}
\label{app:finite-bch-definition}

Represent $\F_{256}$ in the polynomial basis with modulus
$X^8+X^6+X^3+X^2+1$ (hexadecimal \texttt{14d}), and let $\gamma$ be the
element represented by $X$. For $\delta\in\{37,39\}$, let $g_\delta$ be the
binary minimal polynomial product for the roots
$\gamma,\ldots,\gamma^{\delta-1}$. Write $\operatorname{ext}$ for even
parity extension from $255$ to $256$ coordinates. Explicitly, the outer is
\[
 \mathcal B=\operatorname{span}_{\F_2}
 \bigl(\{\operatorname{ext}(X^i g_{39}):0\le i<123\}
 \cup\{\operatorname{ext}(X^j g_{37}):0\le j<5\}\bigr).
\]
The polynomials, with coefficient $i$ stored in bit $i$, are
\begin{center}
\begin{tabular}{@{}ll@{}}
$g_{37}$ & \texttt{1c3b42fe115aa90747020d79ac738adb}\\
$g_{39}$ & \texttt{1f3126ff25af19628f815685cb83da770f}.
\end{tabular}
\end{center}
The quotient $g_{39}/g_{37}$ is primitive irreducible of degree eight.
Thus the five added residues are independent modulo the code generated by
$g_{39}$. The outer has dimension $123+5=128$, is contained in $P$, and
contains $Q$. Even parity and the BCH bound give $d(\mathcal B)\ge38$.
The all-one word belongs to $Q$, hence to $\mathcal B$.

\begin{lemma}[Common spectrum of the intermediate codes]
\leavevmode\par
Every dimension-$128$ code between $Q$ and $P$ has the same weight spectrum.
If $q_w$ counts weight-$w$ words of $Q$ and $h_w$ counts them in one nonzero
$Q$-coset in $P$, then
\begin{equation}
 A_w(P)=q_w+255h_w,\qquad A_w(\mathcal B)=q_w+31h_w.
 \label{eq:finite-quotient-spectrum}
\end{equation}
\end{lemma}
\begin{proof}
A cyclic shift on the punctured coordinates acts on $P/Q$ as multiplication
by $X$ modulo $g_{39}/g_{37}$. Primitivity makes this action transitive on
its $255$ nonzero residues. Parity extension preserves this weight-preserving
action, so all nonzero cosets have the same spectrum. Each intermediate
dimension-$128$ code is a union of $Q$ and $31$ nonzero cosets.
\end{proof}

This lemma connects the implemented generators to the certificate's\linebreak
syndrome-subspace definition of the outer. It does not identify the two
coordinate codes. Independent row permutations make their identical
spectra sufficient for the same first-moment bounds.

\paragraph{Spectrum inequalities.}
Write $A_w:=A_w(\mathcal B)$. Evenness, minimum distance, and complementation
give $A_0=A_{256}=1$, $A_w=A_{256-w}$, and nonzero weights restricted to
\[
 \mathcal W:=\{38,40,\ldots,218\}\cup\{256\}.
\]
The bound inputs comprise individual caps $A_w\le a_w$ and a joint
inequality on the three lowest complementary shell pairs. They are derived
from linear constraints, rather than a model for the unknown spectrum.

The constraints use the exact endpoint spectra of Fujiwara and
Kusaka~\cite{FujiwaraKusaka2021}: the dimension-$71$ extended BCH spectrum
and the dual spectrum of the\linebreak dimension-$187$ extended BCH code.
Generator containments and MacWilliams transforms relate these known
spectra to $P$, $Q$, and their duals. The polynomial
\[
 \mathcal K_j(w):=\sum_v(-1)^v\binom wv\binom{256-w}{j-v}
\]
gives the transform
$A_j(D^\perp)=|D|^{-1}\sum_w A_w(D)\mathcal K_j(w)$ for a binary linear
code $D$. Containment gives coefficientwise inequalities between the
corresponding spectra.

Two additional sources sharpen this relaxation.
First, the hulls $R_P=P\cap P^\perp$ and $R_Q=Q\cap Q^\perp$ have dimensions
$85$ and $93$, with
\[
 R_P\subset R_Q\subset Q,\qquad P\cap R_Q^\perp=Q.
\]
The nonzero cosets of $R_P$ in $R_Q$ also have a common spectrum.
Writing their spectra as $r$ and $s$ gives $A(R_Q)=r+255s$.
For example, the displayed intersection and MacWilliams identity imply
\[
 h\le 2^{-93}\mathcal K(r-s)
\]
coefficientwise. This is a set-difference argument, not subtraction of
unrelated upper bounds. Signed Fourier transforms of
$(-1)^{\wt(c)/2}$ on the even codes give further linear inequalities through
their hulls. Our verification code derives these constraints and checks the requisite
binary ranks and intersections.

Second, a finite Fourier case cover proves $d(Q^\perp)\ge30$.
For a punctured dual word, partition cases by its first nonzero Fourier
cyclotomic coset. In each case, a root run or an independent-set witness
bounds the support size. In the latter witness, Frobenius and coordinate
scaling preserve independence; appending a vector with nonzero coordinate
sum to vectors with zero coordinate sums increases the rank. The final
rank cannot exceed the support size. The exceptional case is split
exhaustively on a further Fourier value. Parity and affine invariance
transfer the resulting bound to the extended dual.

Consequently, projection of $Q$, and of every $Q$-coset, onto any $29$
coordinates is uniform. Otherwise an annihilating linear functional would
give a nonzero dual word of weight at most $29$.
Thus
\[
 \sum_w q_w\mathcal K_j(w)=\sum_w h_w\mathcal K_j(w)=0
 \qquad(1\le j\le29).
\]
These moment equations also hold for $\mathcal B$.
For example, squaring the degree-$14$ reproducing kernel for the binomial
measure gives the cap
\begin{equation}
 A_w\le
 \frac{2^{128}}{\displaystyle\sum_{j=0}^{14}
             \mathcal K_j(w)^2/\binom{256}{j}}.
 \label{eq:finite-christoffel-cap}
\end{equation}
Indeed, normalize that kernel to equal one at $w$; its square is
nonnegative, has degree $28$, and has binomial expectation equal to the
reciprocal denominator. The moment identities transfer that expectation to
a uniform word of $\mathcal B$.

The retained caps take the minimum of such bounds, constant-weight packing
bounds, and exact linear-programming witnesses. The joint relaxation has
$196$ nonnegative variables and $1163$ normalized rows.
Its checker reconstructs the rows and verifies rational dual feasibility
and the objective bound. In particular, if $Fv\le b$, $v\ge0$, and
$y\ge0$ satisfies $F^\top y\ge c$, then $c^\top v\le b^\top y$.
Equality rows have unrestricted multipliers. This exact test, not a solver's
floating feasibility tolerance, supplies the inequality consumed below.
The published endpoint spectra remain mathematical inputs; empirical
spectrum estimates do not.

\subsection{Selected IMT maps and positive transfers}
\label{app:finite-inner-definition}

The half-rate construction uses exactly the maps in Table~\ref{tab:imt-maps}.
In particular, the image of $A$ does not contain the all-one word; the
feedback columns have weight five. No field-multiplier update or identity
$CA=0$ is used. The seven-state fixed-weight transfers
$T_j(z)$ in Appendix~\ref{app:imt-finite-transfers} apply without a limiting
argument. They distinguish zero, arbitrary nonzero mass, and five uniform
image-weight shells. Their cancellation bounds use the feedback map $C$
separately from the expansion $A$.

Fix an output tilt $z\in(0,1)$ and write $E_j(z):=T_j(z)$.
For a shuffled region of $L$ bits, form
\[
 \mathcal E(u,z):=\sum_{j=0}^{128}\binom{128}{j}E_j(z)u^j,\qquad
 R_j(z):=\frac{[u^j]\mathcal E(u,z)^{L/128}}{\binom Lj}.
\]
Conditional on its weight, the region is a uniform slice. The binomial
factors count its stepwise supports, proving this coefficient formula.
Multiplication retains the inner state across steps and regions; the
matrices need not commute. From the zero initial state, the conditional
bad-word probability is at most
$z^{-H}e_ZR_{j_1}\cdots R_{j_{256}}\1$, where $e_Z$ is a row vector.
Truncation above degree $q$ cannot change any coefficient through degree $q$.

For dense input laws the verifier may instead use the complete occupation
or Fourier moment bound of Appendix~\ref{app:imt-finite-transfers}.
Every selection bounds the same terminal moment; an entrywise minimum of
different measure representations would not be justified.

\subsection{One active outer row}

A row of weight $w$ reaches a uniform $w$-subset of the $256$ regions.
Its single position in each reached region is independently uniform.
Use the same $E_0,E_1$ above, including the exact weight-one
cancellation data for the selected independent maps.
Put $\ell=L/t$, $R_0=E_0^\ell$, and
$R_1=\ell^{-1}\sum_{i=0}^{\ell-1}E_0^iE_1E_0^{\ell-1-i}$.
The coefficient bound including the $L$ possible active row positions is
\[
 c_w(z):=\frac{Lz^{-H}}{\binom{256}{w}}\,
 e_Z[u^w](R_0+uR_1)^{256}\1.
\]
Let $Z_{H,q}$ count bad messages with exactly $q$ active rows and put
$F_q:=\E Z_{H,q}$. Each shell may choose its own certified tilt;
write $c_w$ for the resulting upper bound. Then $F_1\le\sum_w c_w A_w$.

Let $J=\{38,40,42\}$. The retained outer witness proves
$\sum_{w\in J}b_w^\star A_w\le B^\star$ for fixed positive rational
$b_w^\star$. Complement symmetry and nonnegativity imply
\begin{equation}
 F_1\le
 B^\star\max_{w\in J}\frac{c_w+c_{256-w}}{b_w^\star}
 +\sum_{w\in\mathcal W\setminus(J\cup(256-J))}c_w a_w.
 \label{eq:finite-q1-domination}
\end{equation}
This reuses a weighted outer inequality, not a probability computed for
another inner. All $c_w$ are evaluated for the selected IMT instance.

\subsection{Sparse occupancies and constant-row compositions}

Partition ordinary weights into twelve bands:
\[
\begin{gathered}
\relax[38,42],\ [44,48],\ [50,54],\ [56,64],\ [66,78],\ [80,100],\\
[102,128],\ [130,154],\ [156,176],\ [178,198],\ [200,210],\ [212,218],
\end{gathered}
\]
retaining only even weights. Keep the all-one word as a thirteenth label.
For an ordinary band $g$, choose $p_g\in(0,1)$ and define
\[
 \Gamma_g=\max_{w\in g}
 \frac{a_w}{\binom{256}{w}p_g^w(1-p_g)^{256-w}},
 \qquad \rho_g\ge\Gamma_g^{1/256}.
\]
The counting measure of a randomly permuted row in band $g$ is pointwise
dominated by $\Gamma_g$ times product Bernoulli$(p_g)$ measure.
Indeed, at a support of weight $w$ its mass is $A_w/\binom{256}{w}$.
The comparison makes different regions independent before their inner states
are propagated. Split $\Gamma_g$ into $256$ equal factors, one per region.

For fixed occupancy $q$, start with $V_j^{(0)}=R_j$ for $0\le j\le q$.
For $0\le r<q$ and $0\le j<q-r$, set
\[
 V_j^{(r+1)}=
 \max_g \rho_g\bigl((1-p_g)V_j^{(r)}+p_gV_{j+1}^{(r)}\bigr).
\]
The maximum is entrywise and includes the exact label $(p_*,\rho_*)=(1,1)$.
Induction dominates every fixed sequence of band choices. There are $13^q$
such sequences, giving
\begin{equation}
 F_q\le\binom Lq\,13^qz^{-H}e_Z(V_0^{(q)})^{256}\1.
 \label{eq:finite-sparse-bound}
\end{equation}
An entrywise maximum may select different bands for different entries.
That only enlarges the envelope; it is not a claim that the actual rows
make those adaptive choices.

For the smaller instances, a sharper composition split keeps $h$ all-one
rows exact and applies the recurrence only to $d$ ordinary rows.
Initialize $V_j^{(0)}=R_{h+j}$ for $0\le j\le d$ and apply $d$ recurrence
steps, maximizing over the twelve ordinary bands.
The composition contribution is bounded by
\[
 U_{d,h}=\binom Ld\binom{L-d}{h}\,12^d
          z^{-H}e_Z(V_0^{(d)})^{256}\1.
\]
For $d=0$ there is no density cost. Summing every $h=0,\ldots,q$ with
$d=q-h$ covers occupancy $q$. Witness choices may vary with the composition.
Only complete, checked sums enter the final ledger.

\subsection{Dense occupancies by a change of input law}

The dense bound compresses the ordinary/all-one compositions into occupancy
and mean reference density. It needs the following pointwise comparison.

\begin{lemma}[Shuffled Bernoulli comparison]
Let $n\ge1$ and let $S$ be a sum of $n$ independent Bernoulli variables with mean $nr$.
For $D_n:=\min\{n+1,4\lceil\sqrt n\rceil\}$ and $0\le j\le n$,
\[
 \Prob[S=j]\le D_n\binom njr^j(1-r)^{n-j}.
\]
The same factor compares the uniformly shuffled bit vector with
$n$ independent Bernoulli$(r)$ bits.
\end{lemma}
\begin{proof}
The logarithm of the moment generating function is bounded by that of
$\operatorname{Bin}(n,r)$, by concavity in each success probability and
Jensen's inequality. Chernoff's bound gives
$\Prob[S=j]\le\exp(-nD_{\rm KL}(j/n\Vert r))$.
The binomial mass on the right without $D_n$ equals this exponential
times the mass $b_{n,j}$ at the mean of $\operatorname{Bin}(n,j/n)$.
That mean is a mode, so $b_{n,j}\ge1/(n+1)$.
The latter binomial has variance at most $n/4$; Chebyshev places at least
$3/4$ of its mass within $\sqrt n$ of its mean. There are at most
$2\sqrt n+1$ integers there, giving
$b_{n,j}\ge3/[4(2\sqrt n+1)]\ge1/(4\sqrt n)$.
The endpoint cases are deterministic. Finally, divide the weight
probabilities by $\binom nj$ to compare uniformly shuffled vectors.
\end{proof}

Fix reference coordinates $\xi_g\in(0,1)$ for the twelve ordinary labels,
and $\xi_*=1$. These coordinates do not change across the density cover.
For an ordered assignment $a$ of $q$ labels, let
$\nu_a=q^{-1}\sum_i\xi_{a_i}$ and $\vartheta_a=q\nu_a/L$.
Choose an input tilt $x>0$, with row probabilities
$p_g=\xi_g/[x+(1-x)\xi_g]$, and define
\[
 C_g(x):=\max_{w\in g}
   \frac{a_wx^w}{\binom{256}{w}\xi_g^w(1-\xi_g)^{256-w}},
 \quad C_*(x)=x^{256},
 \quad B(\vartheta,x)=1-\vartheta+\vartheta/x.
\]
Changing from row probabilities $p_g$ to $\xi_g$ multiplies the density of a
region of weight $j$ by a row-normalization factor and $x^{-j}$.
The shuffle preserves $j$. Apply the lemma, then absorb $x^{-j}$ into
independent Bernoulli inputs of probability
$r=\vartheta/[x+(1-x)\vartheta]$. The product of row costs and normalizers
is exactly $\prod_i C_{a_i}(x)B(\vartheta_a,x)^N$.

For $\lambda>0$, use a valid seven-state IMT transfer to define
\[
 T(r,z):=\sum_{j=0}^t\binom tj r^j(1-r)^{t-j}E_j(z),\qquad
 \mathcal M(r,\lambda):=e^{\lambda H}e_ZT(r,e^{-\lambda})^{N/t}\1.
\]
The fixed-assignment first moment is at most
$D_L^{256}\prod_iC_{a_i}(x)B(\vartheta_a,x)^N\mathcal M(r,\lambda)$.
For an auxiliary real slope $\eta$, put
$S(x,\eta)=\sum_g C_g(x)e^{-\eta\xi_g}$, including the all-one label.
Expanding $S^q$ sums all ordered assignments. Therefore, for a density
interval $I$, let $U_{q,I}$ be the contribution to $\E Z_H$ from
assignments with $q$ active rows and $\nu_a\in I$. Then
\begin{equation}
 \begin{split}
 U_{q,I}\le{}&\binom Lq D_L^{256} S(x,\eta)^q\\
 &{}\cdot\sup_{\nu\in I}
 \left\{e^{\eta q\nu}B(q\nu/L,x)^N
 \mathcal M\!\left(\frac{q\nu/L}{x+(1-x)q\nu/L},\lambda\right)\right\}.
 \end{split}
 \label{eq:finite-dense-tilt}
\end{equation}
Every interval can choose its own witnesses. Fixed coordinates $\xi_g$
ensure that the intervals still cover every assignment.
At $q=L,\nu=1$, the all-one cost cancels the normalization exactly;
no ordinary-band multiplicity is charged to that word.

\paragraph{Finite rectangle certificates.}
The retained checker fixes $r\in(0,1)$ within each rectangle and chooses
$x(\vartheta)=\vartheta(1-r)/(r(1-\vartheta))$.
Since $x$ now depends on $\nu$, it applies the label sum separately at each
distinct $\nu$, and pays the composition count $\binom{q+12}{12}$.
This factor is necessary; the fixed-$x$ bound cannot simply be reused.

For $u=\log x$, the function $\log S(e^u,\eta)$ is convex, because each band
cost is a maximum of exponentials and log-sum-exp preserves convexity.
A secant upper line $a+bu$, with $0\le b\le256$, bounds it on the interval.
After substitution, the scalar exponent becomes at most
\[
 qa+\eta q\nu-qb\log\frac r{1-r}-N\log(1-r)
       +qb\log\vartheta+(N-qb)\log(1-\vartheta).
\]
Both logarithm coefficients are nonnegative.
Replacing the logarithms by tangent upper bounds at an interior
$\vartheta_0$ leaves a function affine in $\nu$ and convex in $q$:
its quadratic coefficient in $q$ is
$b\nu L^{-1}(\vartheta_0^{-1}+(1-\vartheta_0)^{-1})\ge0$.
Thus its maximum occurs at a rectangle corner.
For intervals ending at $\vartheta=1$, use slope $b=256$ and extend by
continuity at the feasible endpoint $q=L,\nu=1$.

A finite partition tree covers every integer $q$ and every reference density
in its declared range. Each leaf includes the number of integer occupancies
it covers, an upper bound on the composition factor, and the fixed-$r$
matrix bound. The exact union sums all leaf bounds.
The original evaluations pay $(L+1)^{256}$; the refined receipts multiply
their bounds by the exact rational factor $(D_L/(L+1))^{256}$.
This changes only the density comparison, not any transfer or cover.

\paragraph{Short-length refinements.}
The two shortest instances need sharper bounds, not a different encoder.
First, activation from zero can retain uniform-shell mass. For a
weight-$j$ input, let $c_j$ bound every nonzero syndrome fiber. A valid
complete zero row retains $E_{ZZ}=k_jz^j/\binom{128}{j}$ and sends
$a_vc_jz^j/\binom{128}{j}$ to the shell of size $a_v$, with no arbitrary
component. Here $k_j$ counts weight-$j$ kernel words and $a_v$ counts
states with image weight $v$. For Bernoulli input, let $b_w$ count
nonzero weight-$w$ words in the image of $C^\top$.
Fourier inversion similarly bounds each
weighted nonzero syndrome by
\[
 h=2^{-19}(1-r+rz)^{128}
 \left(1+\sum_w b_w\left|1-\frac{2rz}{1-r+rz}\right|^w\right).
\]
Keeping the exact zero entry and sending $a_vh$ to each shell is valid.
The verifier compares complete moment bounds, not their individual entries.

Second, the factor $D_L$ can be replaced by a smaller
composition-sensitive bound. Fix an assignment of reference probabilities
$\xi_i$ (zero for inactive rows), with mean $\vartheta$. To compare its
shuffled weight-$j$ probability to $\operatorname{Bin}(L,\vartheta)$,
assume $0<\vartheta<1$ and $0<j<L$, and put
$y=j/L$ and $a=y(1-\vartheta)/(\vartheta(1-y))$.
Exponential tilting gives the exact ratio
\[
 \frac{\prod_i(1-\xi_i+\xi_i a)}
      {(1-\vartheta+\vartheta a)^L}
 \frac{\Prob[\sum_i X_i'=j]}{\Prob[\operatorname{Bin}(L,y)=j]},
 \qquad
 \Prob[X_i'=1]=\frac{\xi_i a}{1-\xi_i+\xi_i a}.
\]
Here the $X_i'$ are independent Bernoulli variables.
Concavity bounds the product in the numerator by
$(1-\nu+\nu a)^q$. If the tilted variance is at least $V$, Fourier
inversion gives the point-mass bound $e^{-V}I_0(V)$, where
$I_0(V):=\pi^{-1}\int_0^\pi e^{V\cos\phi}\,d\phi$.
Indeed, the absolute
characteristic function is at most
$\exp(-2V\sin^2(\phi/2))$, and integration proves the claim.
This bound decreases with $V$.

For an entire $(q,\nu)$ rectangle, minimize each tilted variance
$\xi(1-\xi)a/(1-\xi+\xi a)^2$ over the allowed tilt interval. Its minimum
is at an endpoint. A finite linear program then minimizes the average of
these lower bounds subject to the allowed mean reference density.
Multiplication by the minimum $q$ gives a uniform $V$.
Maximizing the preceding likelihood ratio over integer $j$ gives a
uniform routing factor for the rectangle. Endpoint weights are bounded
directly. This calculation is also valid with a restricted set of labels.

Third, fixing $x$ throughout a rectangle permits the label sum
$S(x,\eta)^q$ in \eqref{eq:finite-dense-tilt} without the
composition-count factor required by the varying-$x$ calculation.
The iid probability $r$ then varies with $\vartheta$.
To keep the moment and normalizer coupled, fix the setup-averaged
nonnegative function $f(u)=\E_\theta[z^{\wt(I_\theta(u))}]$ on $N$-bit inputs
and put $M(r)=\E_{U\sim\mathrm{Bern}(r)^N}f(U)$. For
$D(\vartheta)=1-\vartheta+\vartheta/x$,
\[
 D(\vartheta)^N M\!\left(\frac{\vartheta}{x+(1-x)\vartheta}\right)
 =\sum_u f(u)x^{-\wt(u)}
          \vartheta^{\wt(u)}(1-\vartheta)^{N-\wt(u)}.
\]
Thus, as a function of $v=\log(\vartheta/(1-\vartheta))$,
\[
 F(v)=N\log D(\vartheta)+\log M(r(\vartheta))+\eta L\vartheta
 \quad\hbox{satisfies}\quad F''(v)\ge-(N+|\eta L|)/4.
\]
Indeed, the logarithm of the positive exponential sum is convex,
$(\log(1+e^v))''\le1/4$, and $|\vartheta''(v)|\le1/4$.
The resulting secant bound is
\[
 F(v)\le\max(F(v_0),F(v_1))
           +(N+|\eta L|)(v_1-v_0)^2/32.
\]
Outward endpoint moment bounds can replace $F(v_0),F(v_1)$.
The verifier adds $q\log S+\log\binom Lq+\lambda H$, sums over the
integers $q$ in the leaf, and applies its routing factor.

Finally, split assignments with no high-reference label from those with
at least one, taking a label to be high when $\xi_g>3/4$.
The first part uses $S_{\rm low}$ and the restricted routing comparison.
For the second part, choose $\zeta\le0$, use
$S_{\rm low}+e^{-\zeta}S_{\rm high}$ and multiply by $e^\zeta$.
If $h\ge1$ labels are high, $e^{\zeta h}\le e^\zeta$ proves this upper
bound after expanding the label sum. An analogous split isolates the
all-one label. The two subset bounds are added before comparison with an
unsplit bound. No separately rounded normalizers are subtracted.
Our derivation note \path{SHORT_LENGTH_BOUNDS.md} derives the rectangle
optimizations and records their exact-arithmetic tests.

\subsection{The quarter-rate outer and feedback}
\label{app:finite-quarter-definition}

Over $\F_{128}$ with modulus $X^7+X+1$ and primitive element $X$, let
$p$ and $q$ generate the primitive length-$127$ BCH codes of designed
distances $31$ and $43$. In coefficient-bit notation,
\[
 p=\texttt{d5306d6bfdbc8574719e70d},\qquad
 q=\texttt{5106dae17a61e520c606a4e29}.
\]
The quarter-rate constituent is the span of the parity extensions of
$X^iq$ for $0\le i<29$ and $p,Xp,X^2p$.
The quotient $q/p=\texttt{d5}$ is primitive of degree seven. The same
coset-orbit argument as above proves rank $32$ and spectrum
\[
 W_{32}=W_{29}+\frac7{127}(W_{36}-W_{29}).
\]
Here $W_{29},W_{36}$ are the published extended-BCH spectra retained and
attributed in our author-side
\artifactlink{https://github.com/ladnir/permute_conv/blob/main/workstreams/rate_quarter_bch/SMALLER_OUTER.md}{outer-construction guide}.
The resulting spectrum has minimum weight $32$, multiplicity $588$ at
weight $32$, total mass $2^{32}$, and a single all-one word.

Retain $A$ from Table~\ref{tab:imt-maps} and use $C$ from
Table~\ref{tab:finite-quarter-feedback}. Its columns have weight three.
The fixed-weight and Fourier transfers apply with this map's own dual
spectrum, kernel counts, fiber caps, and weight-one/two cancellation sums.
The quarter-rate calculation uses the exact outer spectrum, so no BCH-256
spectrum inequality is transferred to it.

\begin{table}[tb]
\centering
\caption{Quarter-rate feedback map: bit $i$ of the listed word is
coordinate $i$ of $C^\top(e_j)$.}
\label{tab:finite-quarter-feedback}
\scriptsize
\begin{tabular}{@{}rl@{}}
\toprule
$j$ & $C^\top(e_j)$\\
\midrule
0 & \texttt{808019020040c0800400032500101090}\\
1 & \texttt{01008021009102020030601020308206}\\
2 & \texttt{020006020e0200220400001205680121}\\
3 & \texttt{00100002410a18010108190188020009}\\
4 & \texttt{006040080801003c2000900080604442}\\
5 & \texttt{e1060040101080201088844080801000}\\
6 & \texttt{a0000010020800002064008920020c34}\\
7 & \texttt{4010b19080200090406201800c000000}\\
8 & \texttt{409802c8300001880104400220c40000}\\
9 & \texttt{2808000401442400804120405408a000}\\
10 & \texttt{044140048004110008004c0612100048}\\
11 & \texttt{02202118200840020003264008080210}\\
12 & \texttt{10060025440204005004808802004248}\\
13 & \texttt{00a0aa01990002400290001042002800}\\
14 & \texttt{0001548000c0110c020002284001d800}\\
15 & \texttt{15000c0000240644950a000000020082}\\
16 & \texttt{00000060001128016a000824000121a1}\\
17 & \texttt{0843000024a008508011100001840100}\\
18 & \texttt{1e0c00004200e0010880000011050404}\\
\bottomrule
\end{tabular}
\end{table}

For the two quarter-rate cutoffs, the verifier separately sums all outer
weight classes for $q=1$, band compositions for $q=2$, and region bounds
for $3\le q\le q_{\max}$. It covers the remaining integer types
(zero, ordinary-weight bands, and all-one rows) by a disjoint dense partition.
Here $q_{\max}=64$ for $\delta=0.165$ and $128$ for $\delta=0.19$.
Every dense type uses a product reference law and a complete IMT moment;
the type likelihood ratio is bounded at the vertices of its box.
The union covers $q=1,\ldots,32768$ at each cutoff.
For completeness, the dense comparison can be written directly in terms
of these types. Let $n_g$ count rows of label $g$, so $\sum_g n_g=L$,
and put $J(\boldsymbol n)=L!/\prod_g n_g!$.
Let $\Gamma_g$ dominate the row counting measure by Bernoulli$(p_g)$
bits, as in the sparse calculation. Include zero and all-one labels with
$(\Gamma_0,p_0)=(1,0)$ and $(\Gamma_*,p_*)=(1,1)$.
Choose a positive probability vector $\boldsymbol\pi$ on the labels and
put $r=\sum_g\pi_gp_g$. Independent labels in a region have type
$\boldsymbol n$ with probability $J(\boldsymbol n)\prod_g\pi_g^{n_g}$.
Conditioning on this type gives the permuted reference rows exactly.
Pay that conditioning cost in each of the $B=128$ regions and count
the $J(\boldsymbol n)$ outer assignments once. If $\mathcal U(r,z)$
upper-bounds the complete IMT moment, the contribution of this type is at most
\[
 z^{-H}\mathcal U(r,z)J(\boldsymbol n)^{1-B}
       \prod_g(\Gamma_g\pi_g^{-B})^{n_g}.
\]
Its logarithm is convex in $\boldsymbol n$ for a fixed witness: the only
nonlinear term is $(B-1)\sum_g\log\Gamma(n_g+1)$, which is convex.
Thus vertices of a box intersected with $\sum_g n_g=L$ suffice.
Multiply its maximum by an upper bound on the number of integer types in
the box, then sum the disjoint boxes. This is an auxiliary input-law
comparison, not additional randomness in the encoder setup.
The outward producer, higher-precision replay, and exact-union checker
are separate from the binary64 witness search.

\subsection{Complete finite margins}
\label{app:finite-margin-tables}

Tables~\ref{tab:finite-bch-margins} and~\ref{tab:finite-quarter-margins}
give the complete numerical results for the two finite theorems in
Section~\ref{sec:finite-certificates}. Their margins are certified bounds
on setup failure, not full application-security levels.

\begin{table}[!htbp]
\centering
\caption{Complete finite certificates for the selected $(128,19)$ IMT inner.
Margins are rounded for display; the theorem uses exact rational upper bounds.}
\label{tab:finite-bch-margins}
\begin{tabular}{@{}rrrr@{}}
\toprule
$\log_2K$ & outer rows $L$ & cutoff $H$ & $-\log_2U_K$ (bits) \\
\midrule
16 & 512 & 13\,107 & 41.818327 \\
18 & 2\,048 & 52\,428 & 50.189076 \\
20 & 8\,192 & 209\,715 & 50.062088 \\
22 & 32\,768 & 838\,860 & 48.390492 \\
24 & 131\,072 & 3\,355\,443 & 46.457255 \\
\bottomrule
\end{tabular}
\end{table}

\begin{table}[!htbp]
\centering
\caption{Two distance thresholds for the same quarter-rate encoder at
$K=2^{20}$. The exact unions, not the rounded margins, establish the targets.}
\label{tab:finite-quarter-margins}
\begin{tabular}{@{}rrrr@{}}
\toprule
$\delta$ & cutoff $\lfloor\delta N\rfloor$ & target $\lambda$ & margin (bits) \\
\midrule
0.165 & 692\,060 & 40 & 41.048168 \\
0.190 & 796\,917 & 30 & 30.033491 \\
\bottomrule
\end{tabular}
\end{table}
\FloatBarrier

\subsection{Certificate completion}

The half-rate exact unions have the following component margins, rounded
only for display. Here $q_{\max}$ is the last sparse occupancy.
\begin{center}
\begin{tabular}{@{}rrrrr@{}}
\toprule
$\log_2K$ & $q_{\max}$ & $q=1$ & $2\le q\le q_{\max}$ & $q>q_{\max}$\\
\midrule
16 & 63 & 43.592834 & 60.988751 & 42.317109\\
18 & 127 & 50.189181 & 63.938924 & 79.903280\\
20 & 511 & 50.062272 & 63.000000 & 78.909867\\
22 & 511 & 48.390495 & 67.000000 & 84.999997\\
24 & 511 & 46.457257 & 66.000000 & 191.512160\\
\bottomrule
\end{tabular}
\end{center}
In each case the checked occupancy set is exactly $\{1,\ldots,L\}$.
Summing the three rational bounds and applying
\eqref{eq:finite-certificate} proves \eqref{eq:finite-bch-distance}.
The quarter-rate unions prove Theorem~\ref{thm:finite-quarter-spin}
by the same first-moment inequality.
Injectivity follows independently from sequential inversion of the inner
and injectivity of the block-diagonal outer.

The \artifactlink{https://github.com/ladnir/permute_conv/blob/main/workstreams/inner_design/finite_migration/PAPER_RESULTS.md}{finite IMT results summary}
lists the five half-rate verification receipts, the quarter-rate certificate
and replay, and their exact source paths. Each receipt binds
the selected maps, code parameters, numerical inputs, and verifier sources.

All selected bounds have passed $512$-bit replay, followed by source
authentication, exact coverage checks, and rational summation.
The $K=2^{16}$ dense partition has $1096$ leaves.
Directed arithmetic and upward dyadic conversion supply the numerical
bounds; floating-point search scores do not enter the unions.
Higher-precision replay shares producer formulas, so it is not an
independent proof implementation.

The timing bindings additionally check map identity, outer implementation,
correctness receipts, and measured source and binary identities.
These verification records are separate from the core encoder artifact;
the artifact does not distribute the generated receipts or raw measurements.
This is a computer-assisted proof, not a proof-assistant formalization.
